\documentclass[11pt,letterpaper]{article}
\usepackage[T1]{fontenc}
\usepackage[utf8]{inputenc}
\usepackage{lmodern}
\usepackage[margin=0.85in]{geometry}
\usepackage{tabularx}
\usepackage{booktabs}
\usepackage{amssymb}
\usepackage{enumitem}
\usepackage{xcolor}
\usepackage{array}
\usepackage{parskip}
\usepackage{fancyhdr}

\pagestyle{fancy}
\fancyhf{}
\rfoot{\small Page \thepage}
\lfoot{\small Skill action plan, draft version 2 (for review)}
\renewcommand{\headrulewidth}{0pt}

\newcommand{\bx}{$\square$}

% One item: its text, then the practices the instructor prescribes for it.
\newcommand{\itemhead}[2]{\par\vspace{6pt}\noindent\textbf{#1}\quad #2\par\vspace{-4pt}}
\newenvironment{practices}{\begin{itemize}[leftmargin=2.2em,itemsep=1pt,topsep=2pt,label=\bx]}{\end{itemize}}

\begin{document}

\begin{center}
{\LARGE\bfseries Skill action plan}\\[4pt]
{\large After the start-of-semester survey}\\[2pt]
{\small\color{gray} Draft version 2 for review. A separate screen shown once the survey is submitted; not part of the survey.}
\end{center}

\section*{How it works (notes for review)}

\begin{itemize}[leftmargin=1.5em]
\item \textbf{What the student sees first.} Their average in each of the six blocks, and their three lowest-rated items (ties broken by block average).
\item \textbf{Choose items, not blocks.} The student picks \textbf{one or two items} (recommended: items rather than blocks, since a practice attaches to an item, and ``get better at communication'' is not a plan). The three lowest are suggested; any item may be chosen.
\item \textbf{The instructor supplies the plan.} Each item has two prescribed practices below: a concrete, countable behavior with a target for the semester, in the spirit of exposure practice (do the uncomfortable thing a set number of times). The student picks one practice per item, or writes their own in the same form (action, count, deadline).
\item \textbf{Commitment.} The page records item, practice, and target, and shows them back at the end of the semester with ``Did you do it?'' and ``Did it help?''. Optionally, a mid-semester reminder with the count so far.
\item Targets assume a 14-week semester with a group project, presentations, and data work. Adjust the counts and the named resources (checklists, course modules, events) to the class.
\end{itemize}

\section*{Prescribed practices, by item}

\subsection*{1. Communication}

\itemhead{1.1}{I can present to a class or a group clearly, without reading from my notes.}
\begin{practices}
\item Volunteer to present your group's work; before each, rehearse out loud, standing, twice. Target: 2 presentations.
\item Record yourself answering a two-minute question on your phone, watch it, and redo it once. Target: once a week, 10 weeks.
\end{practices}

\itemhead{1.2}{I can write a clear one-page memo or report that someone else can act on.}
\begin{practices}
\item Write a one-page summary of the week's lecture or reading, topic sentence first. Target: 8 memos.
\item Before every submitted assignment, cut 20\% of the words, then have one classmate read it and say what the main point was. Target: every submission.
\end{practices}

\itemhead{1.3}{In a discussion, I can summarize what someone else said before I respond.}
\begin{practices}
\item In group meetings, restate the previous speaker's point before giving yours, and keep a tally. Target: 3 per meeting.
\item After each class, write the three main points in your own words. Target: 10 classes.
\end{practices}

\subsection*{2. Teamwork and leadership}

\itemhead{2.1}{I do my share of group work on time, without being reminded.}
\begin{practices}
\item Within a day of each group meeting, post your tasks with dates in the group chat. Target: every meeting.
\item Deliver each of your group deliverables two days before the group's internal deadline, and count the misses. Target: zero misses.
\end{practices}

\itemhead{2.2}{When my group disagrees, I can help us settle it without anyone checking out.}
\begin{practices}
\item When the group disagrees, propose a decision rule (vote, try both, ask the professor) before taking a side. Target: 3 occasions.
\item Read the one-page conflict guide (instructor supplies) and apply one technique at the next disagreement, then note what happened. Target: 2 occasions.
\end{practices}

\itemhead{2.3}{When a group has no plan, I can set one up: tasks, owners, and deadlines.}
\begin{practices}
\item At the first project meeting, be the one who writes the task-owner-deadline table, and update it weekly. Target: all semester.
\item Run a meeting with an agenda you sent the day before. Target: 3 meetings.
\end{practices}

\subsection*{3. Critical thinking and data}

\itemhead{3.1}{Facing a messy question, I can break it into smaller parts I can answer.}
\begin{practices}
\item For each problem set, write the sub-questions before solving anything. Target: 6 problem sets.
\item For the project, draw the question as a tree (claim, parts, evidence needed for each) at each milestone. Target: 3 trees.
\end{practices}

\itemhead{3.2}{I can use numbers or data to support or reject a claim.}
\begin{practices}
\item Each week, find one statistic in the news and check it against the original source. Target: 10 checks.
\item In the project, before analyzing, write one claim and the number that would prove it wrong. Target: every milestone.
\end{practices}

\itemhead{3.3}{I can judge whether a source, a statistic, or an AI answer is trustworthy.}
\begin{practices}
\item For every source in the project, write one line: who produced it, how, and what they gain. Target: every source.
\item Once a week, ask an AI tool a question you know the answer to and find its error. Target: 8 weeks.
\end{practices}

\subsection*{4. Technology}

\itemhead{4.1}{I can use spreadsheet or statistical software (Excel, R, Stata) to analyze data.}
\begin{practices}
\item Redo one class example in Excel or R each week without looking at the solution. Target: 8 examples.
\item Complete the free course module (instructor names it) with its exercises. Target: 4 exercises.
\end{practices}

\itemhead{4.2}{I can teach myself a new software tool from its documentation or videos.}
\begin{practices}
\item Pick a tool you do not know (pivot tables, Git, Zotero) and learn it from the documentation to do one real task. Target: 2 tools.
\item When stuck, spend 15 minutes with the documentation before asking anyone, and count the times. Target: 10 times.
\end{practices}

\itemhead{4.3}{I can use AI tools for schoolwork in a way I could openly explain to a professor or employer.}
\begin{practices}
\item For one assignment, keep an AI log: prompts, what you used, what you changed, and attach it. Target: 1 assignment.
\item Add a three-line disclosure of AI use to every submission. Target: every submission.
\end{practices}

\subsection*{5. Professionalism}

\itemhead{5.1}{I meet deadlines and commitments without last-minute excuses.}
\begin{practices}
\item In the first week, put every deadline for all your courses in one calendar with a reminder three days before. Target: done by week 1.
\item Submit every assignment at least 24 hours early, and count the misses. Target: zero misses.
\end{practices}

\itemhead{5.2}{I check my work for errors before I submit it.}
\begin{practices}
\item Use the five-item checklist (instructor supplies) before every submission and tick it. Target: every submission.
\item Read each submission aloud once before sending. Target: 8 submissions.
\end{practices}

\itemhead{5.3}{I communicate with professors and employers professionally (email, meetings, follow-up).}
\begin{practices}
\item Send every email to a professor or employer from the template (subject, greeting, one ask, signature). Target: every email.
\item Attend office hours with a written question, and send a two-line thank-you afterwards. Target: 3 visits.
\end{practices}

\subsection*{6. Career and self-development}

\itemhead{6.1}{I can name my two strongest and two weakest skills for the job I want.}
\begin{practices}
\item Ask three people who know your work to name one strength and one weakness, and write them down. Target: 3 people.
\item Take the free assessment (instructor names it) and write a half-page reaction. Target: 1 assessment.
\end{practices}

\itemhead{6.2}{I know what the job I want requires and what I still lack.}
\begin{practices}
\item Read ten postings for the job you want, list the five most common requirements, and mark which you have. Target: 10 postings.
\item Do one informational interview with someone in that job. Target: 1 interview.
\end{practices}

\itemhead{6.3}{I reach out to people I do not know (alumni, professors, professionals) to learn about careers.}
\begin{practices}
\item Contact two new people a month (alumni on LinkedIn, professors, professionals) with a three-line message asking one question. Target: 8 contacts.
\item Attend two career-center or department events and introduce yourself to two people at each. Target: 4 introductions.
\end{practices}

\section*{What the student fills in}

\begin{tabularx}{\linewidth}{@{}lX@{}}
\toprule
Item & \rule{8cm}{0.4pt} \\
Practice & \bx\ first \quad \bx\ second \quad \bx\ my own: \rule{5cm}{0.4pt} \\
Target & \rule{3cm}{0.4pt} by \rule{3cm}{0.4pt} \\
\bottomrule
\end{tabularx}

\vspace{4pt}
(A second block appears if the student chose two items.)

\end{document}
